\documentclass[11pt,a4paper]{article}

% ---- packages (all standard; compiles on Overleaf as-is) ----
\usepackage[utf8]{inputenc}
\usepackage[T1]{fontenc}
\usepackage{lmodern}
\usepackage{amsmath}
\usepackage[margin=2.2cm]{geometry}
\usepackage{xcolor}
\usepackage{booktabs}
\usepackage{longtable}
\usepackage{array}
\usepackage{enumitem}
\usepackage{parskip}
\usepackage[hidelinks]{hyperref}

\definecolor{tvblue}{HTML}{2563EB}
\definecolor{tvgrey}{HTML}{444444}
\newcommand{\term}[1]{\textcolor{tvblue}{\textbf{#1}}}
\newcommand{\q}[1]{\medskip\noindent\textbf{#1}\par\nobreak}
\setlist{itemsep=2pt, topsep=3pt}
\newcolumntype{L}[1]{>{\raggedright\arraybackslash}p{#1}}

\title{\textbf{Teraval --- Internal Team Report (Detailed)}\\[4pt]
\large A complete, plain-English walkthrough of the whole project}
\author{Team: Kartik Joshi, Prem Kukreja, Gagandeep Singh, Samuel Alex, Aditya Chitale\\
Module: Corporate Finance \quad$\cdot$\quad Instructor: Dr.\ Nathaniel Christopher}
\date{Repository: \texttt{github.com/KartikJoshi23/TeravalAI}\\
Live: \texttt{teraval-ai.vercel.app}}

\begin{document}
\maketitle

\begin{quote}
\small \textbf{How to read this document.} This is our own team notebook, written so \emph{anyone}
on the team --- finance or not --- can understand every part of the project and defend it in the
viva. It follows the assignment brief section by section (1--11) and answers all five project
questions in full. Where a big word appears, it is defined on the spot. Formulas are shown and then
explained in plain English, with our real numbers plugged in. Nothing is left as ``it just works.''
\end{quote}

\tableofcontents
\newpage

% =====================================================================
\section{Introduction (Brief \S1)}

\subsection{The project topic}
We chose \term{Topic 9 --- AI Capital-Budgeting Dashboard}. \term{Capital budgeting} is the part of
corporate finance that decides whether a big, long-term investment is worth making: you estimate all
the future cash the project will bring in and take out, pull those future amounts back to
``today's money,'' and check whether what you get back is worth more than what you put in. Our
dashboard does exactly this for one concrete, realistic project, and layers artificial-intelligence
tools on top to forecast, stress-test and explain the decision. (Topic~9 naturally absorbs Topic~10,
``relevant cash flows,'' because our build-vs-rent analysis is a textbook relevant-cash-flow problem.)

\subsection{The company and the decision-maker}
The company is \term{Barq AI}, a (hypothetical but realistic) UAE artificial-intelligence
infrastructure operator. ``Barq'' means \emph{lightning} in Arabic --- fitting for a business that
sells raw computing speed. The decision-maker is Barq's \term{CFO} and \term{investment committee}:
the people who must sign off on spending several billion dirhams. The dashboard is the tool they
would take into that meeting.

\subsection{The financial problem}
Barq must decide: \emph{should we spend about \textbf{AED 5.8 billion} to build and own a
\textbf{40~MW} AI/GPU data-centre hall in Abu Dhabi, and operate it for \textbf{8 years}, renting out
its computing power to AI companies?} A \term{data centre} is a building full of powerful computers;
a \term{GPU} (graphics processing unit) is the chip that does the heavy maths for AI, rented out here
by the GPU-hour. 40~MW is the electrical power the computers draw --- a very large facility.

\subsection{Why the decision is important}
\begin{itemize}
  \item \textbf{Bet-the-company scale.} AED 5.8~billion is enormous; getting it wrong could sink
  the firm.
  \item \textbf{It is largely irreversible.} Purpose-built power and liquid cooling cannot easily be
  repurposed if the project fails --- you cannot un-pour the concrete.
  \item \textbf{The revenue price is volatile.} The price to rent a GPU for one hour fell from about
  \$8 to roughly \$2.85--3.50 as supply grew. If the price stays low, the project can lose money ---
  this single number decides almost everything.
  \item \textbf{Timing and location help.} The UAE is deliberately becoming a global AI-computing hub,
  and Abu Dhabi's electricity is cheap, which lowers the biggest running cost.
\end{itemize}

\subsection{The purpose of the dashboard}
To replace gut feeling with numbers: to compute every standard capital-budgeting metric correctly and
live, to show them as clear visuals, to stress-test the decision against bad luck, to add genuine AI
(forecasting and a trained risk model) that improves the analysis, and to end with a defensible
recommendation --- while always leaving the final call to a human.

% =====================================================================
\section{Corporate Finance Concepts, with formulas (Brief \S2)}
These are the finance ideas the project uses. Each has its formula and a plain-English meaning.

\subsection{Time value of money and discounting}
Money today is worth more than the same money later, because today's money can be invested. To
compare cash arriving in different years, we shrink each future amount to its \term{present value}
(PV):
\[
  \text{PV} \;=\; \frac{\text{FV}}{(1+r)^{n}},
\]
where FV is the future amount, $r$ is the yearly discount rate, and $n$ is the number of years away.
Example: AED~100 in 5 years at $r=9\%$ is worth $100/1.09^{5}\approx\text{AED }65$ today.

\subsection{Relevant (incremental) cash flows}
Only cash flows that \emph{change because of the decision} count. Sunk costs (already spent) and
unaffected costs are ignored. This is why our build-vs-rent test compares Barq's project
\emph{against the next-best way of getting the same computing} (renting), not against ``doing
nothing'' --- Barq needs the compute either way.

\subsection{Discounted cash flow (DCF) and Net Present Value (NPV)}
\term{DCF} means valuing something by discounting all its future cash flows. The headline DCF metric
is \term{NPV}: add up every year's cash flow, each discounted to today, including the up-front spend
(a negative number at year~0):
\[
  \text{NPV} \;=\; \sum_{t=0}^{n} \frac{\text{CF}_t}{(1+r)^{t}}.
\]
\textbf{Rule: NPV $>0$ means the project creates value --- accept; NPV $<0$ means it destroys value
--- reject.} It is the single most important number in the project.

\subsection{Internal Rate of Return (IRR) and Modified IRR (MIRR)}
The \term{IRR} is the discount rate at which NPV would be exactly zero --- effectively the project's
own yearly \% return:
\[
  \text{find } r^{*} \text{ such that } \sum_{t=0}^{n}\frac{\text{CF}_t}{(1+r^{*})^{t}} = 0,
  \qquad \text{IRR}=r^{*}.
\]
If IRR $>$ WACC, the project beats its cost of money. IRR has a known flaw (it assumes cash is
reinvested at the IRR itself), so we also compute \term{MIRR}, which reinvests positive cash at the
WACC instead:
\[
  \text{MIRR} \;=\; \left(\frac{\text{FV of positive flows at WACC}}{-\,\text{PV of negative flows}}\right)^{1/n} - 1 .
\]

\subsection{Profitability Index (PI)}
Value created per dirham invested:
\[
  \text{PI} \;=\; \frac{\text{PV of the inflows}}{\text{initial outlay}}
  \;=\; 1 + \frac{\text{NPV}}{\lvert \text{CF}_0 \rvert}.
\]
PI $>1$ means every AED~1 invested returns more than AED~1 of value --- accept.

\subsection{Payback and discounted payback}
\term{Payback} is how many years until the cumulative cash flow first turns positive (you have your
money back). \term{Discounted payback} does the same but on \emph{discounted} cash flows, so it is
always longer and more honest. Short payback = less risk exposure.

\subsection{The discount rate: WACC}
The rate $r$ we discount with is the \term{Weighted Average Cost of Capital} --- the blended yearly
return the company's investors (debt + equity) require:
\[
  \text{WACC} \;=\; \frac{E}{V}\,r_e \;+\; \frac{D}{V}\,r_d\,(1-T),
\]
where $E,D$ are equity and debt, $V=E+D$, $r_e,r_d$ are the costs of equity and debt, and $T$ is the
tax rate (interest is tax-deductible, hence the $(1-T)$). We build ours up from the official Central
Bank of the UAE risk-free anchor (base rate 3.65\%, 3-month EIBOR 3.94\%), a levered equity risk
premium, and an after-tax cost of debt --- which lands near 8.6\%. \textbf{We deliberately kept 9\%},
a slightly tougher bar, so our ``accept'' is conservative, not flattering.

\subsection{Depreciation and the tax shield}
Big assets are \term{depreciated} (their cost is spread over their useful life for accounting). We
use \term{straight-line} depreciation (equal amount each year). Depreciation is not cash, but it
reduces taxable profit, so it creates a \term{tax shield} (cash saved on tax) that we add back inside
operating cash flow.

\subsection{Equivalent Annual Cost (EAC) and the annuity factor}
To compare two options that last \emph{different numbers of years} (own the hall for 8 years vs. a
3-year cloud rental), you cannot compare their totals fairly. Instead you convert each into a level
yearly cost with the same present value --- its \term{EAC}:
\[
  \text{EAC} \;=\; \frac{\text{PV of all the option's costs}}{a(r,n)},
  \qquad
  a(r,n) \;=\; \frac{1-(1+r)^{-n}}{r},
\]
where $a(r,n)$ is the \term{annuity factor} (the present value of AED~1 received each year for $n$
years). Lower EAC wins. (An \term{annuity} is a fixed payment each period; a \term{perpetuity} is one
that never ends --- we use the annuity idea directly in EAC.)

\subsection{Break-even}
The \term{break-even} value of an input is the value at which NPV is exactly zero. Our key one is the
break-even GPU rental price: \textbf{\$3.34/GPU-hr}. Below it, the project loses money. We find it by
\term{bisection} (repeatedly halving a search range until NPV crosses zero).

\subsection{Monte-Carlo and probability of loss}
Instead of one guess per uncertain input, we give each a \emph{range} and draw thousands of random
combinations, computing NPV for each. The share of runs with NPV $<0$ is the \term{probability of
loss}. This turns ``it might go wrong'' into a number ($\approx$21\%).

% =====================================================================
\section{Data and Assumptions (Brief \S3)}
\label{sec:data}
Every input is tagged \term{Historical (H)}, \term{Current (C)}, \term{Forecast (F)},
\term{User-entered (U)} or \term{AI-generated (AI)}. Money is in AED millions; USD converts at the
pegged \$1 = AED~3.6725.

\begin{center}
\small
\begin{longtable}{L{4.4cm} L{4.2cm} L{6cm}}
\toprule
\textbf{Input} & \textbf{Base value} & \textbf{Type and basis}\\
\midrule
\endhead
GPU rental price & \$4.00 / GPU-hr & \textbf{U/C} --- assumed contracted rate; market is \$2.85--3.50\\
Utilization & 80\% & \textbf{U} --- mix of contracted + on-demand demand\\
Electricity tariff & AED 0.15 / kWh & \textbf{C} --- Abu Dhabi industrial (official figure pending)\\
PUE (cooling overhead) & 1.20 & \textbf{C} --- typical for liquid-cooled GB200\\
WACC (discount rate) & 9\% & \textbf{F} --- built from official CBUAE rates + risk premia\\
Facility scale & 40 MW; 300 racks; 21{,}600 GPUs & \textbf{C} --- GB200 NVL72 platform\\
Rack cost & \$3.5 m / rack & \textbf{C} --- vendor list pricing\\
Facility build cost & \$13 m / MW & \textbf{C} --- AI-grade data-centre build\\
Project life & 8 years & \textbf{U} --- appraisal horizon (GPU useful life)\\
Tax rate & 9\% & \textbf{C} --- UAE corporate tax (from 2023)\\
Working capital & AED 150 m & \textbf{U} --- recovered at the end\\
CBUAE base rate & 3.65\% & \textbf{H/C} --- Central Bank of the UAE, 29 Jul 2026 (evidence file)\\
3-month EIBOR & 3.94\% & \textbf{H/C} --- CBUAE EIBOR table, 31 Jul 2026 (evidence file)\\
GPU-price history & 40-month series, \$5.6$\to$\$3.1 & \textbf{H} --- representative 2023--26 market trend (trains the ML forecaster)\\
Forecast spot rate & $\approx$\$2.73 / GPU-hr & \textbf{AI} --- output of the trained AR(1) forecaster\\
Accept-probability & 95\% at base & \textbf{AI} --- output of the trained risk classifier\\
\bottomrule
\end{longtable}
\end{center}

\textbf{What each data type means here:}
\begin{itemize}
  \item \textbf{Historical} --- the past: the 40-month GPU-price series and the dated CBUAE rate
  screenshots in the \texttt{datasets/} folder.
  \item \textbf{Current} --- today's market/structural facts (hardware costs, tariff, tax, platform).
  \item \textbf{Forecast} --- forward-looking judgement, chiefly the 9\% WACC.
  \item \textbf{User-entered} --- the six dashboard sliders an analyst moves (price, utilization,
  WACC, tariff, PUE, rack cost) plus the horizon.
  \item \textbf{AI-generated} --- \emph{genuinely} produced by our trained models: the forecast spot
  price and the accept-probability (see Section~\ref{sec:ml}). This is real, not a relabelled formula.
\end{itemize}

\textbf{Official-data note.} We downloaded the CBUAE base rate and EIBOR (with evidence files). The
Abu Dhabi industrial electricity tariff could not be downloaded --- the \texttt{data.bayanat.ae}
portal was down --- so we kept the AED~0.15/kWh benchmark and documented the gap honestly.

% =====================================================================
\section{Financial Calculations, worked (Brief \S4)}
The brief asks for at least three major calculations; we have many. Here is the base case built up
from scratch.

\subsection{Step 1 --- the up-front spend (Year 0)}
\begin{itemize}
  \item IT hardware: $300\text{ racks}\times\$3.5\text{m}\times3.6725 = \text{AED }3{,}856\text{ m}$.
  \item Facility: $40\text{ MW}\times\$13\text{m}\times3.6725 = \text{AED }1{,}910\text{ m}$.
  \item Capex $=3{,}856+1{,}910=\text{AED }5{,}766\text{ m}$; plus working capital AED~150~m
  $\Rightarrow$ \textbf{initial outlay AED 5{,}916 m} (this is $\text{CF}_0$, negative).
\end{itemize}

\subsection{Step 2 --- the yearly operating cash}
\begin{itemize}
  \item Revenue $=21{,}600\text{ GPUs}\times8{,}760\text{ hrs}\times80\%\times\$4.00\times3.6725
  \approx \textbf{AED 2{,}224 m/yr}$.
  \item Energy $\approx$ AED~57~m; total operating cost $\approx$ \textbf{AED 600 m/yr}; so
  \textbf{EBITDA $\approx$ AED 1{,}624 m/yr} (earnings before interest, tax, depreciation).
  \item Subtract depreciation, apply 9\% tax on the profit, add the tax shield back $\Rightarrow$
  operating cash flow. A GPU \term{refresh} outflow hits in Year~4; salvage + working-capital
  recovery come in Year~8.
\end{itemize}

\subsection{Step 3 --- discount and score}
Discounting the 8-year free-cash-flow stream at 9\% gives the headline results:

\begin{center}
\begin{tabular}{lll}
\toprule
\textbf{Metric} & \textbf{Value} & \textbf{What it says} \\
\midrule
NPV & \textbf{+AED 1{,}854 m} & positive $\rightarrow$ creates value $\rightarrow$ accept \\
IRR & 16.5\% & beats the 9\% WACC $\rightarrow$ accept \\
MIRR & 12.6\% & still beats 9\% \\
Profitability Index & 1.31 & $>1 \rightarrow$ AED 1.31 of value per AED spent \\
Payback & 5.0 years & within the 8-year life \\
Discounted payback & 6.3 years & within life, but late \\
Break-even GPU price & \textbf{\$3.34/GPU-hr} & inside today's \$2.85--3.50 market --- thin! \\
\bottomrule
\end{tabular}
\end{center}

\subsection{Step 4 --- the relevant-cash-flow test: build vs rent}
Because Barq needs the compute either way, we compare \emph{building} against \emph{renting} the same
capacity from a hyperscaler, on \term{EAC} (different lives). Building is a fixed cost; renting scales
with usage. Result: \textbf{building wins only above about 77\% sustained utilization}; at the base
case building beats renting by a thin \textbf{incremental NPV of +AED 320 m}. As a self-check, the
sign of that incremental NPV matches the difference in EACs --- two independent routes agree.

\subsection{Step 5 --- the decision indicated}
Every deterministic metric points the same way at the base case: \textbf{accept}. But the break-even
price sits right inside the current market and the safety margin is thin --- which is exactly why the
scenario, sensitivity, Monte-Carlo and AI sections exist: to test whether ``accept'' survives bad luck.

% =====================================================================
\section{AI Features --- the five-point breakdown (Brief \S5)}
The brief wants at least five AI features, each explained five ways. We have \textbf{seven}, and
\textbf{two of them are genuinely trained machine-learning models} (Section~\ref{sec:ml}), not
rule-based helpers. Every feature is \emph{grounded}: the deterministic engine computes all
valuations; the AI forecasts inputs, quantifies risk, or explains --- it never invents a valuation.

\newcommand{\feat}[6]{%
  \medskip\noindent\textbf{#1}\par
  \begin{itemize}[leftmargin=1.4em,itemsep=1pt,topsep=1pt]
    \item \textbf{What it does:} #2
    \item \textbf{Information it uses:} #3
    \item \textbf{Result it produces:} #4
    \item \textbf{How it helps the decision-maker:} #5
    \item \textbf{Limitation / risk:} #6
  \end{itemize}}

\feat{F1 --- Trained GPU-price forecaster (machine learning)}
{Learns how the GPU rental price moves over time and projects it forward with a confidence band.}
{A 40-month historical price series; it fits an AR(1) time-series model by least squares.}
{A forward price path with P10--P90 bands and a held-out test error (RMSE $\approx$\$0.06, MAPE 1.7\%);
central forecast $\approx$\$2.73/GPU-hr.}
{Replaces a flat guess with a data-driven, validated estimate of the most important and most volatile
input, and flags that spot prices sit below break-even.}
{The market is non-stationary and the training series is representative, not a licensed feed --- so it
is shown as a band, never as certainty.}

\feat{F2 --- Trained risk classifier / surrogate (machine learning)}
{Learns the accept/reject boundary of the whole model, so it can instantly predict the chance a given
set of assumptions creates value.}
{$\sim$4{,}000 driver combinations, each labelled accept/reject \emph{by our own deterministic
engine}; a logistic-regression classifier is trained on a 75/25 train/test split.}
{A live accept-probability, a learned price$\times$utilization ``decision surface,'' and permutation
feature-importance --- with held-out test accuracy 96.7\% and ROC-AUC 0.998.}
{Gives an instant risk read as sliders move and, via feature-importance, an independent, data-driven
confirmation that price and utilization dominate.}
{It emulates the current model, so it inherits the model's own assumptions --- it is not an external
opinion.}

\feat{F3 --- Monte-Carlo risk simulator}
{Runs thousands of random ``what-if'' worlds to measure the spread of possible outcomes.}
{Triangular ranges around each key input; 5{,}000 draws through the engine with a fixed random seed.}
{The probability of loss ($\approx$21\%), the average NPV (+AED 1{,}456 m), and P10/P90 outcomes.}
{Turns vague uncertainty into a concrete risk number the committee can weigh.}
{Results depend on the assumed input ranges; wrong ranges give a falsely tight or loose picture.}

\feat{F4 --- AI Finance Assistant (NVIDIA NIM chatbot)}
{Answers natural-language questions about the dashboard in plain English.}
{The live computed model state for the current tab, handed to an NVIDIA NIM language model.}
{A written, tab-aware explanation streamed word by word.}
{Lets a non-financial manager interrogate the analysis conversationally.}
{Language models can \emph{hallucinate}; we mitigate by grounding it in the engine's numbers and
telling it to explain only, never generate figures --- and it falls back to a local answerer if the
server is off.}

\feat{F5 --- Scenario generator}
{Auto-builds a coherent optimistic / base / pessimistic set of assumptions.}
{The current slider values, shifted by sensible amounts in the favourable/unfavourable directions.}
{Three ready-made scenarios you can apply with one click.}
{Saves manual set-up and keeps the three cases internally consistent.}
{The shift sizes are judgement calls, not laws of nature.}

\feat{F6 --- Anomaly / unrealistic-assumption detector}
{Flags inputs that are outside believable ranges.}
{The current assumptions vs. physical/market limits (e.g.\ a PUE below~1 is impossible).}
{A warning when an input is implausible.}
{Stops ``garbage-in, garbage-out'' before it reaches the decision.}
{It checks plausibility, not correctness --- a believable-but-wrong number still passes.}

\feat{F7 --- Recommendation engine}
{Writes the final accept/reject paragraph and a stage-gate plan.}
{NPV, IRR, PI, break-even, the dominant driver, and the probability of loss.}
{A short, grounded conclusion with conditions attached when the ``accept'' is thin.}
{Gives the committee a ready, defensible narrative tied to the numbers.}
{It synthesises; it does not decide --- the human must still sign.}

% =====================================================================
\section{How the machine learning actually works (train / test)}
\label{sec:ml}
\textbf{Important correction.} An earlier version of this project had \emph{no} machine learning ---
its ``AI'' was simulations plus a chatbot --- and we honestly said so. That has now changed: we added
two genuinely trained models (F1 and F2 above) that \emph{do} use a train/test split and report
held-out metrics. Here is exactly how, so we can defend it.

\subsection{What a train/test split is, and why it matters}
In machine learning you \term{train} a model on part of the data and \term{test} it on a separate part
it never saw during training. If it also does well on the unseen \emph{test} data, it has genuinely
\emph{learned the pattern} rather than memorised the answers (\term{overfitting}). Reporting the
\emph{test} score --- not the training score --- is what makes the claim honest.

\subsection{The risk classifier (F2), step by step}
\begin{enumerate}
  \item \textbf{Make data.} We draw $\sim$4{,}000 random combinations of the five drivers and ask the
  \emph{deterministic engine} for the truth of each: value-creating (NPV~$>0$) or not. So the labels
  are our own finance model --- the AI cannot fabricate; it can only learn what the engine already
  says.
  \item \textbf{Split.} 75\% train, 25\% test (3{,}000 / 1{,}000), shuffled with a fixed seed.
  \item \textbf{Train.} A \term{logistic regression} --- $p = 1/(1+e^{-(\mathbf{w}\cdot\mathbf{x}+b)})$
  --- is fitted by gradient descent to minimise cross-entropy loss. We add a price$\times$utilization
  feature because revenue depends on their product.
  \item \textbf{Test.} On the held-out 1{,}000: \textbf{accuracy 96.7\%, ROC-AUC 0.998, F1 0.968}
  (precision 0.973, recall 0.963). Train accuracy (98.4\%) is close to test accuracy, so it is not
  overfit.
  \item \textbf{Use.} It outputs a live accept-probability (95\% at base) and, by \term{permutation
  importance} (scramble one input, see how much accuracy falls), ranks the drivers:
  price$\times$util $>$ price $>$ utilization $>$ WACC $>$ tariff $>$ PUE --- which independently
  matches our tornado chart.
\end{enumerate}

\subsection{The price forecaster (F1), step by step}
It fits a \term{first-order autoregression} --- each month depends on the previous one:
\[
  p_t \;=\; c + \phi\,p_{t-1} + \varepsilon_t, \qquad
  \text{long-run level} = \frac{c}{1-\phi}.
\]
Fitted by \term{ordinary least squares} on the first $\sim$29 months and tested on the last~10, it
gives $\phi=0.955$ (a stable, mean-reverting process), a long-run level of \$2.63, and a held-out
one-step \textbf{RMSE of \$0.063 (MAPE 1.7\%)}. Forecasting forward 8 years gives an average spot
price of $\approx$\$2.73 (P10--P90 \$2.27--3.19).

\subsection{The honest, useful conclusion}
The forecaster tracks the \emph{spot/market} rate; our base case uses \$4.00, a \emph{contracted}
rate. The gap between them (\$2.73 spot vs \$4.00 contracted, with break-even at \$3.34) is exactly
the ``contracted premium'' Barq must secure. So the AI does not contradict the project --- it
\emph{quantifies why the recommendation insists on locking in multi-year customer contracts} rather
than riding the falling spot market. Both models \emph{augment} the analysis; the deterministic engine
still computes every valuation.

% =====================================================================
\section{The dashboard, component by component (Brief \S5 dashboard)}
The app has \textbf{five tabs} (grouped so nothing is buried), a header (logo, live self-test badge,
``Download summary'' button, current verdict), and a floating AI assistant on every tab. The brief's
six required components are all present (KPI cards, cash-flow chart, scenario comparison, sensitivity,
risk panel, AI recommendation) plus several extras.

\begin{enumerate}
  \item \textbf{Overview.} Opens with a \term{Decision Brief} that states the question, the stakes and
  \emph{how to read the page} --- so the verdict lands as a conclusion, not a surprise. Then the live
  \term{3D data-centre hall} (racks brighten with utilization; the heat plume warms as cooling gets
  less efficient; click a rack to inspect), the six \term{KPI cards} (NPV, IRR, MIRR, PI, payback,
  break-even --- each labelled accept/reject), and the \term{AI recommendation}.
  \item \textbf{Cash Flow \& Build-vs-Rent.} The 8-year cash-flow chart (Year~0 the big spend, Year~4
  the refresh dip, a line crossing zero at the discounted payback), then the Build-vs-Rent / EAC
  comparison.
  \item \textbf{Scenarios \& Sensitivity.} The three scenarios side by side; \term{decision
  thresholds} (the exact point the verdict flips); six live \term{sliders}; a \term{tornado chart}
  ranking each driver's impact; and a \term{risk-and-alert panel}.
  \item \textbf{AI \& Forecasting.} The Monte-Carlo distribution, the scenario generator, the rate
  forecast band --- and the \term{Predictive-AI} section showing the two trained ML models (metrics,
  confusion matrix, decision surface, feature-importance, live accept-probability).
  \item \textbf{Board \& Ethics.} A \term{Departmental Review Board} (five departments each score the
  same live case by a fixed rule and a weighted verdict aggregates them --- base case: ``approve with
  conditions''), then the ethics panel, the data-provenance panel, and the assumptions-audit table
  with a live model self-test.
\end{enumerate}
Every component reads one shared ``memory box'' of assumptions, so moving any slider updates the whole
app at once. \textbf{Download summary} prints a clean one-page PDF of the current live case.

% =====================================================================
\section{The AI Finance Assistant --- sample questions and answers (Brief \S6)}
The floating assistant answers questions about the dashboard, grounded in the live numbers. Five (plus)
worked examples:

\begin{description}[leftmargin=0pt,style=nextline]
  \item[Q1. ``What is the project's NPV, and should we accept it?''] The base-case NPV is
  \textbf{+AED 1{,}854 m}, so it creates value and the signal is \emph{accept} --- but conditionally,
  because the safety margin is thin.
  \item[Q2. ``Why does the NPV fall when the GPU price drops?''] Revenue is directly proportional to
  the GPU price, so a lower price lowers every year's cash flow; below \$3.34/GPU-hr the NPV turns
  negative.
  \item[Q3. ``Which assumption has the greatest effect on the result?''] The GPU rental price (rank~1
  on the tornado and in the ML feature-importance), followed by utilization.
  \item[Q4. ``What happens if utilization falls to 65\%?''] The project moves toward the pessimistic
  case and can turn negative; renting also becomes cheaper than building below $\sim$77\% utilization.
  \item[Q5. ``Should management accept or reject the project?''] Accept, \emph{as a stage-gated
  commitment}: release capital in tranches tied to measured utilization, lock in multi-year contracts,
  and keep a rent fallback.
  \item[Q6. ``Explain the result to a non-financial manager.''] ``We put in about AED~5.8~bn; over
  8~years the discounted cash we get back is about AED~1.85~bn more than that, a 16.5\% return --- good,
  but it depends heavily on keeping GPU rental prices up, so we should sign customers first.''
  \item[Q7. ``What is the probability of losing money?''] About \textbf{21\%}, from a 5{,}000-run
  Monte-Carlo simulation.
\end{description}
If the AI server is off, a built-in, engine-grounded answerer gives the same factual replies.

% =====================================================================
\section{Scenario and Sensitivity Analysis (Brief \S7)}

\subsection{The three scenarios}
\begin{center}
\begin{tabular}{L{2.4cm} L{5.6cm} l l}
\toprule
\textbf{Scenario} & \textbf{Key assumptions} & \textbf{NPV} & \textbf{Decision}\\
\midrule
Optimistic & \$6.0/hr, 90\% util, cheaper power, 8\% WACC & +AED 10{,}216 m & Accept\\
Base & \$4.0/hr, 80\% util, 9\% WACC & +AED 1{,}854 m & Accept (conditional)\\
Pessimistic & \$2.5/hr, 65\% util, dearer power, 11\% WACC & $-$AED 4{,}138 m & Reject\\
\bottomrule
\end{tabular}
\end{center}

\subsection{The four things the brief asks us to explain}
\begin{itemize}
  \item \textbf{Which variable has the greatest impact?} The \textbf{GPU rental price} --- the tornado
  chart and the ML feature-importance agree. Utilization is second.
  \item \textbf{Under what conditions does the decision change?} It flips to reject when the price falls
  below the \$3.34 break-even, or utilization drops under $\sim$67\%, or the two combine (the
  pessimistic case). The decision-thresholds panel shows each flip point live.
  \item \textbf{The main financial risk.} A structural, lasting fall in the GPU rental rate below
  break-even --- and the Monte-Carlo says there is a $\approx$21\% chance the base case loses money.
  \item \textbf{The best management response.} De-risk the revenue: sign multi-year contracted offtake
  above break-even, stage the capital release, and keep renting available as a fallback below the 77\%
  build-vs-rent crossover.
\end{itemize}

% =====================================================================
\section{Ethical Use of AI (Brief \S8)}
Each ethics point the brief lists, with our real safeguard:
\begin{itemize}
  \item \textbf{Accuracy of AI estimates.} Forecasts are shown as \emph{ranges} (P10--P90), and the ML
  models report \emph{held-out test error} so their accuracy is stated, not assumed.
  \item \textbf{Risk of incorrect data.} Every input is source-tagged and range-checked; the anomaly
  detector flags implausible values.
  \item \textbf{AI hallucinations.} The chatbot is \emph{grounded} --- handed the computed numbers and
  told to explain only --- so it cannot invent figures.
  \item \textbf{Confidentiality.} All computation is client-side; commercially sensitive assumptions
  never leave the browser.
  \item \textbf{Bias in forecasts.} A bull-market bias is explicitly stress-tested by the pessimistic
  scenario and the Monte-Carlo downside.
  \item \textbf{Human review.} A live self-test checks the maths, and every recommendation states that
  the AI advises only.
  \item \textbf{Responsibility.} \textbf{The final invest/reject decision rests with the CFO and the
  investment committee, not with the AI.}
\end{itemize}

% =====================================================================
\section{Final Recommendation (Brief \S9)}
\textbf{Proceed with the investment --- conditionally, as a stage-gated commitment.} The base case
creates about AED~1.85~billion of value at a 16.5\% return against a conservative 9\% hurdle, and the
build-vs-rent and board analyses agree. But the margin is thin: break-even (\$3.34) sits close to the
\$4.00 price, the data-driven forecast puts the \emph{spot} rate below break-even, and there is a
$\approx$21\% modelled chance of loss. So management should: \textbf{(i)} secure multi-year contracted
offtake to hold the effective rate above break-even; \textbf{(ii)} release the capital in tranches
gated on measured utilization staying above roughly 67--77\%; and \textbf{(iii)} keep hyperscaler
rental as a fallback. This conclusion is supported by the financial calculations, the dashboard, the
scenario and sensitivity analysis, the AI-generated insights (forecaster + risk classifier), and our
own judgement. The internal Board Review reaches the same verdict: \emph{approve with conditions}.

% =====================================================================
\section{The Five Project Questions, answered (Brief \S11)}

\q{Q1. What financial decision is being analysed, and why is it important?}
Whether Barq AI should commit $\approx$AED~5.8~bn to build and own a 40~MW Abu Dhabi GPU data centre
and operate it for 8 years. It matters because it is bet-the-company scale, largely irreversible, and
its payoff hinges on a single volatile price (GPU rental), in a strategically timed UAE AI-hub market.

\q{Q2. Which finance concepts, formulas, data and assumptions are required?}
Time value of money and discounting; DCF/NPV; IRR/MIRR; PI; payback and discounted payback; WACC;
depreciation and the tax shield; relevant (incremental) cash flows; EAC and the annuity factor;
break-even; and Monte-Carlo (all with formulas in Section~2). Data: capex build-up, revenue/opex,
the six drivers, official CBUAE rates, and a historical price series --- classified H/C/F/U/AI in
Section~3.

\q{Q3. What do the calculations show across the three scenarios?}
Optimistic +AED~10{,}216~m (accept), base +AED~1{,}854~m (accept, conditional), pessimistic
$-$AED~4{,}138~m (reject). The decision genuinely flips with the assumptions, which is the whole point
of the appraisal.

\q{Q4. How can AI improve the analysis, dashboard, forecasting and decision-making?}
Four ways, all present: a \emph{trained forecaster} turns the key volatile input into a validated,
banded estimate; a \emph{trained risk classifier} gives an instant accept-probability and a
data-driven driver ranking; \emph{Monte-Carlo} converts uncertainty into a probability of loss; and a
\emph{grounded chatbot} explains everything in plain language --- while the deterministic engine keeps
every number honest.

\q{Q5. What final recommendation, and what risks?}
Accept, conditionally and stage-gated (Section~11). The dominant \emph{financial} risk is a structural
fall in the GPU rental price below break-even; the dominant \emph{AI-related} risk is over-trusting an
optimistic forecast --- mitigated by ranges, held-out testing, grounding, and mandatory human review.

% =====================================================================
\section{Technology stack, and how we keep it correct}
\textbf{Front-end:} React~19 + TypeScript, Vite, TailwindCSS~v4, Zustand (shared state), Framer
Motion + GSAP (animation), Recharts (2D charts), Three.js + react-three-fiber (the 3D hall). \textbf{ML
layer:} pure TypeScript (\texttt{web/src/lib/ml/}) --- logistic regression, AR(1), metrics --- so the
maths is inspectable, seeded and testable, with no heavy black-box library. \textbf{Finance engine:}
a TypeScript module mirrored exactly by a Python reference (\texttt{docs/finance-model-reference.py});
the two must always agree. \textbf{Assistant back-end:} FastAPI + NVIDIA~NIM. \textbf{Quality:}
\textbf{45 automatic tests} (Vitest), a \texttt{tsc} type-check and an \texttt{oxlint} pass, a live
in-app self-test of financial identities, and grounding so the AI never invents numbers.
\textbf{Live:} front-end on Vercel, AI back-end on Render.

% =====================================================================
\section{Mini-glossary}
\begin{description}[leftmargin=0pt, style=nextline, itemsep=1pt]
  \item[Capital budgeting] deciding whether a big long-term investment is worth it.
  \item[GPU / GPU-hour] an AI chip / renting one for an hour (our unit of revenue).
  \item[PUE] Power Usage Effectiveness --- cooling overhead; 1.0 is perfect, 1.2 is good.
  \item[WACC] the blended yearly return investors require; our discount rate (9\%).
  \item[NPV / IRR / MIRR / PI] the core value and return scores (see Section~2).
  \item[EAC / annuity factor] tools to fairly compare options of different lengths.
  \item[Deterministic] same inputs always give the same output (our engine).
  \item[Monte-Carlo / seed] thousands of random what-if runs / a fixed start so they repeat.
  \item[Machine learning] a model that \emph{learns} a pattern from example data.
  \item[Train/test split] learn on part of the data, score on unseen part (honesty check).
  \item[Logistic regression] a model that predicts a probability of a yes/no outcome.
  \item[AR(1)] a time-series model where each step depends on the previous one.
  \item[Accuracy / ROC-AUC / RMSE / MAPE] how right a model is (classification / forecast error).
  \item[Feature importance] how much each input matters to a trained model.
  \item[Surrogate model] a fast learned stand-in for a slower exact model.
  \item[Grounded AI] an AI that may only use the real numbers it is handed.
\end{description}

\vspace{1em}
\begin{center}\small\textcolor{tvgrey}{%
Teraval $\cdot$ Barq AI capital-budgeting appraisal $\cdot$ every figure verified against the
deterministic reference model; the two ML models are reported on held-out test data.}\end{center}

\end{document}
